\section{Primitive role architecture}\label{sec:primitive-role-architecture}

% Mathematical job:
%   - Fix the six role meanings P1 through P6 at the repaired alignments.
%   - State local boundary distinction rules (P1/P5, P1/P4, P1/P2, P1/P3,
%     P1/P6, P2/P5, ...): distinctions under scoped bookkeeping, not global
%     independence.
%   - Introduce the common behavioral / verification / structural-categorical
%     level pattern and the selected-lift atlas vocabulary.
%   - Use Foundations (Tsiokos 2026) terminology: theory package
%     T = (Z, f, Sigma_f, E, A), closure operator, idempotent endomap.
%
% Repaired role meanings (non-negotiable):
%   P1  update-vs-quotient compatibility / descent obstruction
%   P2  macro-specification representability
%   P3  enriched route/coherence mismatch
%   P4  stage/order/transition/refinement structure
%   P5  quotient packaging / quotienting
%   P6  audit/bookkeeping/provenance/replay/checkability
%
% Forbidden regressions:
%   generic transition assigned to P3
%   generic path/derivation assigned to P5
%   generic semantics/interpretation assigned to P1
%
% Governing sources:
%   workspace/meta/lens_lab/synthesis/six_primitive_state_table.md
%   workspace/meta/lens_lab/synthesis/six_primitive_boundary_matrix.md
%   workspace/meta/lens_lab/synthesis/six_primitive_level_lift_atlas.md
%   workspace/meta/lens_lab/synthesis/six_primitive_global_nonclaims.md
%   docs/papers/Tsiokos_2026_Six_Birds_Foundations_of_Emergence_Calculus.tex
%     (executive P1-P6 map; finite theory package; three-certificate loop)
%   paper/prior_literature_context.md (Tier A1)
%
% Overclaim risks:
%   - Restating Paper 1 as a full six-primitives algebra.
%   - Losing the repaired role meanings under casual prose.
%   - Treating boundary rules as global independence theorems.

This section fixes the meanings of the six typed closure roles
\Pone{}, \ldots, \Psix{}, the behavioral, verification, and
structural-categorical level pattern they share, the local boundary
distinctions between them, and the selected-lift atlas vocabulary.
These meanings are fixed for the remainder of the paper and govern
every subsequent definition, lower bound, and upper-bound
classification.

\subsection{Role meanings}\label{subsec:role-meanings}

\begin{definition}[Primitive closure roles]\label{def:primitive-roles}
The six primitive closure roles are
\begin{itemize}[topsep=2pt,itemsep=1pt]
  \item \Pone{} --- update-vs-quotient compatibility / descent
    obstruction~\citep{GrothendieckSGA1,JanelidzeTholen1994};
  \item \Ptwo{} --- macro-specification
    representability~\citep{MacLane1998};
  \item \Pthree{} --- enriched route/coherence
    mismatch~\citep{BaaderNipkow1998,Nakahara2003}; \Pthree{} is not a
    directionality certificate;
  \item \Pfour{} --- stage/order/transition/refinement
    structure~\citep{Milner1989,BackVonWright1998};
  \item \Pfive{} --- quotient packaging /
    quotienting~\citep{KemenySnell1960};
  \item \Psix{} --- audit/bookkeeping/provenance/replay/%
    checkability~\citep{GreenKarvounarakisTannen2007,Merkle1987}.
\end{itemize}
\end{definition}

\begin{remark}[Corrected alignments]\label{rem:repaired-alignments}
Three alignments that Definition~\ref{def:primitive-roles} rules out
are used throughout later classification: generic transition data is
not \Pthree{}; generic path or derivation data is not \Pfive{}; and
generic semantic or interpretation data is not \Pone{}. These
constraints determine which feature clusters can later be interpreted
as active projections of each role.
\end{remark}

\begin{remark}[Drive face of \Psix{}]\label{rem:p6-drive}
Arguments about affinities, detailed balance, or honest directionality
use the narrower \(\mathrm{P6}_{\mathrm{drive}}\) specialization of
\Psix{}: its affinity-carrying, detailed-balance-breaking
face~\citep{Seifert2012}. \(\mathrm{P6}_{\mathrm{drive}}\) is a face
of \Psix{}, not a replacement for it. A \Pthree{}
route/coherence-mismatch witness does not establish directionality
without a \Psix{}-side audit.
\end{remark}

\subsection{Behavioral, verification, and structural-categorical levels}\label{subsec:levels}

Each role admits a common three-level pattern.

\begin{definition}[Level trichotomy]\label{def:level-trichotomy}
For each role \Pone{}, \ldots, \Psix{}, we distinguish
\begin{itemize}[topsep=2pt,itemsep=1pt]
  \item the \emph{behavioral level}: status, occurrence, or observed
    behavior of the role in a description;
  \item the \emph{verification level}: proof, witness, certificate, or
    checkable record supporting a behavioral
    claim~\citep{MartinLof1984};
  \item the \emph{structural-categorical level}: reusable structure
    supporting the primitive role.
\end{itemize}
Forgetting proceeds from structural-categorical to verification to
behavioral and is lossy. Lifting proceeds from behavioral to
verification to structural-categorical only after adding explicitly
recorded data, and selected lifts are not unique.
\end{definition}

\begin{remark}[Level content per primitive]\label{rem:per-primitive-levels}
The level trichotomy instantiates as follows.
\begin{itemize}[topsep=2pt,itemsep=1pt]
  \item \Pone{}: descent success or descent obstruction status; descent
    proof or obstruction witness; quotient data equipped with a descended
    map, or an explicit no-descended-map obstruction record.
  \item \Ptwo{}: representability or non-representability status;
    representability or non-representability certificate;
    macro-specification layer together with a realization relation.
  \item \Pthree{}: witness behavior of route disagreement; generated route
    terms and non-reducibility certificates; categorical route grammar or a
    reduced categorical universe.
  \item \Pfour{}: stage status, transition occurrence, or refinement
    behavior; order/preorder, transition, or refinement certificates;
    stage/refinement architecture.
  \item \Pfive{}: packaging or identification behavior; equivalence and
    quotient-validity certificate; quotient object with a packaging map
    and universal-property or tower/coarsening data.
  \item \Psix{}: audit-relevant event occurrence or record--event
    association; checkable audit record; audit schema with provenance,
    replay, and check relations.
\end{itemize}
The meta-theoretic properties of this trichotomy --- lossy forgetting,
selected non-unique lifts, and the absence of lower-to-higher determinacy
without added data --- are formalized in
Section~\ref{sec:admissibility-meta-theory}.
\end{remark}

\subsection{Local boundary distinctions}\label{subsec:boundary-matrix}

With role meanings and level pattern fixed, we turn to how the six
roles remain distinct. Table~\ref{tab:boundary-matrix} records the
pairwise local rules. These are local distinctions under scoped
bookkeeping, not a global primitive-independence theorem: two roles
may interact legitimately in a single description, but neither
collapses into the other. The table is reproduced in landscape
orientation on the page immediately preceding the references.

\begin{proposition}[Boundary distinctions under scoped bookkeeping]\label{prop:boundary-distinctions}
For every pair of roles listed in Table~\ref{tab:boundary-matrix}, the
corresponding ``forbidden collapse'' is ruled out as a locally admissible
identification of those two roles under the repaired alignments of
Remark~\ref{rem:repaired-alignments} together with the honest bookkeeping
required by a \FATCD{}.
\end{proposition}

\begin{remark}[Not a global independence theorem]\label{rem:not-global-independence}
The boundary matrix asserts local distinctions, not global
independence. Primitives interact legitimately in a single
description: \Pone{} descent is audited by \Psix{}; \Pthree{} route
terms may be indexed by \Pfour{} stages; nontrivial \Ptwo{}
representability is typically carried over \Pfive{} quotient structure.
No claim is made that the six roles are independent generators of a
global algebra. The typed-non-collapse counterpart of this observation
appears in Section~\ref{sec:admissibility-meta-theory}.
\end{remark}

\subsection{Selected-lift atlas and global nonclaims}\label{subsec:lift-atlas}

The remaining architectural data are the selected lifts between levels
and the nonclaims that accompany them. Upgrading from behavioral or
verification content to the structural-categorical level requires
explicitly recorded added data. The resulting lifts are selected, not
unique, and the atlas carries structural content that is lost under
forgetting.

\begin{definition}[Selected-lift atlas]\label{def:selected-lift-atlas}
The \emph{selected-lift atlas} for the six roles records, for each
\Pone{}, \ldots, \Psix{}:
\begin{itemize}[topsep=2pt,itemsep=1pt]
  \item \Pone{}: proof or witness together with quotient structure, carrying
    either a descended map or an explicit no-map obstruction;
  \item \Ptwo{}: representability or non-representability certificates
    together with the macro-specification layer and realization relation;
  \item \Pthree{}: the selected lift data recorded in the promoted atlas,
    carrying generated route terms, non-reducibility certificates, and
    categorical route grammar or a reduced categorical universe;
  \item \Pfour{}: validity certificates together with structural
    stage/refinement architecture;
  \item \Pfive{}: quotient-validity certificates together with the
    structural quotient data, a packaging map, and universal-property or
    tower/coarsening data;
  \item \Psix{}: record-event or check evidence together with the audit
    schema and its provenance and replay/check relations.
\end{itemize}
Each lift is selected and not unique, and forgetting from a higher level
loses structural content that is not recoverable without re-adding data.
\end{definition}

\begin{remark}[Global nonclaims for the role architecture]\label{rem:atlas-global-nonclaims}
The role architecture preserves the following nonclaims, all carried
forward into the admissibility regime of
Section~\ref{sec:admissibility-meta-theory} and the theorem chain of
Sections~\ref{sec:lower-bounds}--\ref{sec:decomposition-exhaustion}:
\begin{itemize}[topsep=2pt,itemsep=1pt]
  \item no full six-primitives algebra is claimed;
  \item no global primitive-independence theorem is claimed;
  \item no empirical realization is claimed;
  \item no uniqueness of primitive presentations is claimed;
  \item no uniqueness of selected lifts is claimed;
  \item no lower-level-to-higher-level determinacy is claimed;
  \item no instrument-free formulation is claimed;
  \item direct cross-level comparison remains inadmissible without
    translation.
\end{itemize}
\end{remark}
